\documentclass[11pt]{article}
\usepackage[margin=1in]{geometry}
\usepackage[T1]{fontenc}
\usepackage{lmodern}
\usepackage{microtype}
\usepackage{xcolor}
\usepackage{hyperref}
\usepackage{booktabs}
\usepackage{enumitem}
\usepackage{listings}
\usepackage{tcolorbox}
\usepackage{parskip}
\usepackage{array}
\newcolumntype{L}[1]{>{\raggedright\arraybackslash}p{#1}}

% \path{} (from url, loaded by hyperref) prints paths and identifiers
% verbatim -- so `--` is not turned into a dash -- and may break at / _ - .
\makeatletter
\g@addto@macro\UrlBreaks{\do\_\do\-\do\.}
\makeatother
\urlstyle{tt}
\setlength{\emergencystretch}{3em}

\hypersetup{
  colorlinks=true,
  linkcolor=blue!60!black,
  urlcolor=blue!60!black
}

\definecolor{codebg}{RGB}{245,245,245}
\definecolor{codecomment}{RGB}{90,100,90}
\definecolor{codekeyword}{RGB}{0,70,140}
\definecolor{codestring}{RGB}{150,60,20}

\lstdefinestyle{pythonstyle}{
  language=Python,
  backgroundcolor=\color{codebg},
  basicstyle=\ttfamily\small,
  keywordstyle=\color{codekeyword}\bfseries,
  commentstyle=\color{codecomment}\itshape,
  stringstyle=\color{codestring},
  showstringspaces=false,
  breaklines=true,
  frame=single,
  rulecolor=\color{black!20},
  columns=fullflexible,
  keepspaces=true
}
\lstdefinestyle{shellstyle}{
  language=bash,
  backgroundcolor=\color{codebg},
  basicstyle=\ttfamily\small,
  showstringspaces=false,
  breaklines=true,
  frame=single,
  rulecolor=\color{black!20},
  columns=fullflexible,
  keepspaces=true
}
\lstset{style=pythonstyle}

\newtcolorbox{keyidea}{
  colback=blue!5, colframe=blue!50!black,
  title=Key Idea, fonttitle=\bfseries
}
\newtcolorbox{tipbox}{
  colback=green!5, colframe=green!45!black,
  title=Tip, fonttitle=\bfseries
}
\newtcolorbox{warningbox}{
  colback=orange!7, colframe=orange!60!black,
  title=Common Pitfall, fonttitle=\bfseries
}

\title{
  \textbf{Codebase Change Summary --- Week 7}\\[4pt]
  \large Coverage-Based Testing\\[2pt]
  \normalsize SE 413/513 --- Software Testing
}
\author{Demo Testing Suite}
\date{}

\begin{document}
\maketitle

\begin{center}
\textit{What changed in \path{demo-testing-suite/} since Week 6, and why}
\end{center}

\vspace{0.3cm}

% ============================================================
\section{Purpose of This Document}
% ============================================================
The \path{demo-testing-suite/} codebase is a single, shared project that
grows one lecture at a time rather than being rebuilt from scratch each
week. This document summarizes exactly what was added or changed for
Week 7 and, more importantly, explains \emph{why} each change exists ---
which concept from the Week 7 reading, \textit{Coverage-Based Testing:
Measuring and Interpreting Coverage with \texttt{coverage.py} and
\texttt{pytest-cov}}, it demonstrates. Read this alongside the reading
and \path{demo-testing-suite/README.md}.

\begin{keyidea}
Nothing from Weeks 1--6 was removed. This week does three things: it
adds a small, self-contained project that reproduces every coverage
report in the reading; it applies the reading's method to the
\emph{real} suite, closing the two coverage gaps that have been open
since Week 4; and it adds a second small project showing how
\texttt{coverage.py} and \texttt{pytest-cov} relate. The two intentional \texttt{xfail} results and one
intentional \texttt{FAIL} are unchanged.
\end{keyidea}

% ============================================================
\section{What Changed At a Glance}
% ============================================================
\begin{center}
\small
\begin{tabular}{@{}L{5.9cm}L{1.5cm}L{7.0cm}@{}}
\toprule
\textbf{File} & \textbf{Status} & \textbf{Purpose} \\
\midrule
\multicolumn{3}{@{}l}{\textit{The main suite: closing real gaps}}\\
\path{tests/orders/test_order_service.py} & Modified & A declined-payment
stub: \path{place_order()}'s \texttt{"Payment failed"} line had never
run. \\
\path{tests/notifications/test_password_reset_service.py} & Modified &
Two \path{SmtpEmailClient} tests against a patched \texttt{smtplib.SMTP}
--- no mail server needed. \\
\path{app/notifications/password_reset.py} & Modified & One justified
\texttt{\#~pragma:\ no~cover} (an abstract method); docstrings updated. \\
\midrule
\multicolumn{3}{@{}l}{\textit{The reading, reproduced (not in the default run)}}\\
\path{demos/coverage_reports/checkout_app/checkout.py} & New & The
reading's running example, verbatim: every line number matches. \\
\path{demos/coverage_reports/tests/} & New & The v1, v2,
coupons-only, and no-assertion suites from the reading. \\
\path{demos/coverage_reports/coverage_settings.toml} & New & The
reading's \texttt{[tool.coverage.*]} configuration. \\
\path{demos/coverage_reports/README.md} & New & Every command, in the
reading's order, with its expected result. \\
\midrule
\multicolumn{3}{@{}l}{\textit{\texttt{coverage.py} vs.\ \texttt{pytest-cov} (not in the default run)}}\\
\path{demos/coverage_vs_pytest_cov/tempconv/temps.py} & New & A tiny
temperature module: the code under test. \\
\path{demos/coverage_vs_pytest_cov/tests/test_temps.py} & New & A
deliberately thin suite (67\% branch). \\
\path{demos/coverage_vs_pytest_cov/report.py} & New & A command-line
script pytest never runs, measured with \texttt{coverage run}. \\
\path{demos/coverage_vs_pytest_cov/coverage_settings.toml},
\texttt{README.md} & New & Shared settings, and every command with its
expected output. \\
\midrule
\multicolumn{3}{@{}l}{\textit{Tooling and documentation}}\\
\texttt{requirements.txt} & Modified & Pins \texttt{coverage>=7.5} (the
HTML report's Functions view); adds \texttt{pytest-xdist} (parallel
test workers). \\
\texttt{.gitignore} & Modified & Ignores \texttt{coverage.xml} and
\texttt{coverage.json}. \\
\path{docs/DEMO_GUIDE_week07_coverage_based_testing.md} & New &
Section-by-section lecture companion for the reading. \\
\texttt{README.md} (suite and repository) & Modified & Updated layout,
coverage commands, expected output, and the Weekly Updates index. \\
\bottomrule
\end{tabular}
\end{center}

\begin{tipbox}
Reinstall the requirements before running anything this week:
\begin{lstlisting}[style=shellstyle]
cd demo-testing-suite
pip install -r requirements.txt
\end{lstlisting}
\end{tipbox}

% ============================================================
\section{The Reading's Reports, Reproduced}
% ============================================================
\textbf{Location:} \path{demos/coverage_reports/}

This folder is a miniature project laid out exactly like the reading's
(\path{checkout_app/checkout.py}, \path{tests/test_checkout_v1.py},
\ldots), so that report paths and line numbers match the reading
character for character. \textbf{Run its commands from inside the
folder:}
\begin{lstlisting}[style=shellstyle]
cd demos/coverage_reports
\end{lstlisting}

\begin{center}
\footnotesize
\begin{tabular}{@{}L{5.2cm}L{4.4cm}L{5.0cm}@{}}
\toprule
\textbf{Run} & \textbf{Result} & \textbf{Reading section} \\
\midrule
v1, \texttt{coverage run} / \texttt{report} & \texttt{TOTAL 81\%} (test
files included) & The Basic Report \\
v1, \texttt{-{}-source=checkout\_app}, \texttt{-m} & 74\%; \texttt{8,
16-17, 24, 28} & Showing What Was Missed \\
v1, \texttt{-{}-branch} & 66\%; \texttt{BrPart 3} & Turning On Branch
Coverage \\
v1, settings file (\texttt{-{}-rcfile}) & exit 2:
\texttt{66 < fail-under=90} & Making Settings Permanent \\
v1, \texttt{pytest -{}-cov-fail-under=80} & 3 passed, exit 1 & Enforcing
a Minimum \\
coupons-only & 41\%; \texttt{16->18} & Reading the \texttt{Missing}
Column \\
v2 & 100\%, 9 passed & Closing the Gaps \\
no-assertion test & 100\%, 1 passed & Limitation 1 \\
v2 vs.\ \texttt{amount > 50} (live edit) & 100\%, 9 passed & Limitation
3 \\
\bottomrule
\end{tabular}
\end{center}

\begin{warningbox}
Do not add a docstring to \path{checkout_app/checkout.py} or reformat
it. Every \texttt{Missing} entry in the reading names a line of this
file, so moving a line breaks the whole demo.
\end{warningbox}

\subsection*{Why \texttt{coverage\_settings.toml}, not \texttt{pyproject.toml}}
The reading's configuration turns on branch measurement, restricts
\texttt{source}, and sets \texttt{fail\_under = 90}. If the folder had a
\texttt{pyproject.toml} containing it, \emph{every} run would use those
settings, and you could never see the 81\% $\rightarrow$ 74\%
$\rightarrow$ 66\% progression. So the settings live under a name
nothing reads by default, and the ``permanent settings'' step names the
file explicitly:
\begin{lstlisting}[style=shellstyle]
coverage run --rcfile=coverage_settings.toml -m pytest tests/test_checkout_v1.py
coverage report --rcfile=coverage_settings.toml
\end{lstlisting}
\begin{lstlisting}[style=shellstyle]
Coverage failure: total of 66 is less than fail-under=90
\end{lstlisting}
\texttt{pytest -{}-cov -{}-cov-config=coverage\_settings.toml} reads the
same file. In your own project, call it \texttt{pyproject.toml} and
drop the flag.

% ============================================================
\section{From Report to Tests, on the Real Suite}
% ============================================================
\textbf{Locations:} \path{tests/orders/test_order_service.py},
\path{tests/notifications/test_password_reset_service.py},
\path{app/notifications/password_reset.py}

Since Week 5, the whole-suite report has shown two gaps:
\begin{lstlisting}[style=shellstyle]
app/notifications/password_reset.py      19      7      0      0    63%   34, 47, 50-56
app/orders/order_service.py               8      1      2      1    80%   37
\end{lstlisting}
Following the reading's ``From Report to Tests'' method, each gap was
read as a \emph{behavior} that had never been checked, not as a line
number to hit:

\begin{center}
\small
\begin{tabular}{@{}L{3.7cm}L{5.0cm}L{5.7cm}@{}}
\toprule
\textbf{Gap} & \textbf{Behavior never checked} & \textbf{What was done} \\
\midrule
\path{order_service.py},\newline line 37 & A declined payment is reported as
\texttt{"Payment failed"} & New \path{DecliningPaymentStub}: a second
canned answer next to Week 4's \path{PaymentStub} \\
\path{password_reset.py},\newline lines 47, 50--56 & \path{SmtpEmailClient} sends to
the configured host, and closes the connection \emph{even when sending
fails} (the \texttt{finally}) & Two tests that patch
\texttt{smtplib.SMTP} (Week 5's mocking), so no mail server is
involved \\
\path{password_reset.py},\newline line 34 & None: it is the abstract
\path{EmailClient.send}, which only raises
\texttt{NotImplementedError} & Excluded with
\texttt{\#~pragma:\ no~cover} and a comment giving the reason \\
\bottomrule
\end{tabular}
\end{center}

\begin{lstlisting}[style=shellstyle]
pytest --cov=app --cov-branch --cov-report=term-missing:skip-covered
\end{lstlisting}
\begin{lstlisting}[style=shellstyle,language={}]
TOTAL     113      0     34      0   100%

22 files skipped due to complete coverage.
\end{lstlisting}

\begin{tipbox}
Compare the two decisions in \path{password_reset.py}. The abstract
method is excluded because it has no behavior, and that reason fits in
one comment that would pass code review. \path{SmtpEmailClient} was
\emph{not} excluded, even though it looked hard to test: being hard to
test is not a reason to exclude code (Best Practice 7). A test double
made it easy after all.
\end{tipbox}

% ============================================================
\section{What 100\% Still Does Not Tell You}
% ============================================================
Every \path{app/} module is now at 100\% branch coverage. The demo
folder shows why that is not the same as ``fully tested'':
\begin{itemize}
\item \path{tests/test_checkout_noassert.py} reaches 100\% with
\textbf{no assertions} (Limitation 1).
\item In the v2 suite, \texttt{amount >= 50} is never true; the only
True outcome of line 22 comes from \texttt{is\_member} (Limitation 2).
\item Changing \texttt{>=} to \texttt{>} on line 22 leaves v2 passing,
at 100\% (Limitation 3). Make the edit, run, then restore with
\texttt{git checkout -{}- checkout\_app/checkout.py}.
\end{itemize}

\begin{warningbox}
The same is true of the main suite. Week 1's \path{calculate_late_fee}
bug and Week 3's shipping boundary bug both sit inside files that report
100\%. Those defects are found by the \texttt{xfail} tests, not by the
coverage number.
\end{warningbox}

% ============================================================
\section{\texttt{coverage.py} vs.\ \texttt{pytest-cov}: How They Relate}
% ============================================================
\textbf{Location:} \path{demos/coverage_vs_pytest_cov/}

The reading uses both tools, and they produced identical numbers. This
second mini-project shows \emph{why}: \texttt{pytest-cov} is a
\texttt{pytest} plugin that drives the \texttt{coverage.py} engine, so the
two share one data file (\texttt{.coverage}) and one configuration. It
also shows when each is the better choice. It uses its own
\path{tempconv} package, so its line numbers are independent of
\path{checkout_app}. As before, \textbf{run its commands from inside the
folder}:
\begin{lstlisting}[style=shellstyle]
cd demos/coverage_vs_pytest_cov
\end{lstlisting}

\begin{center}
\footnotesize
\begin{tabular}{@{}L{6.4cm}L{3.6cm}L{4.6cm}@{}}
\toprule
\textbf{Run} & \textbf{Result} & \textbf{What it shows} \\
\midrule
\texttt{pytest -{}-cov=tempconv -{}-cov-branch}, and
\texttt{coverage run -{}-branch -{}-source=tempconv -m pytest}
& both 67\%; \texttt{7, 10} & Same engine, same numbers \\
\texttt{pytest -{}-cov ...}, then \texttt{coverage report -m} & same
67\% table, no re-run & One shared \texttt{.coverage} file \\
\texttt{coverage run ... report.py -5 30} & 83\%; \texttt{9} & Measuring
a script, not a test run \\
then \texttt{coverage run -{}-append ... -m pytest} & 100\% & Combining
two runs \\
\texttt{pytest -n 2 -{}-cov=tempconv -{}-cov-branch} & 67\% &
\texttt{pytest-cov} measures the xdist workers \\
\texttt{coverage run ... -m pytest -n 2} & 0\% & Bare
\texttt{coverage run} sees only the main process \\
\texttt{coverage run -m pytest} (no \texttt{-{}-source}) & \texttt{TOTAL 85\%}
& Test files inflate the total \\
\bottomrule
\end{tabular}
\end{center}

\begin{keyidea}
Use \texttt{pytest-cov} for test runs, especially in CI or with parallel
workers. Use \texttt{coverage.py} directly to measure something that is
not a \texttt{pytest} run (a script, a manual run), to combine data from
several runs, or to produce more reports from data you already have.
Either way, it is the same measurement.
\end{keyidea}

% ============================================================
\section{Updated: Expected Test Results}
% ============================================================
From the project root:
\begin{lstlisting}[style=shellstyle]
cd demo-testing-suite
pip install -r requirements.txt
pytest -v
\end{lstlisting}
you should now see \textbf{72 passed, 2 xfailed, 1 failed}: three more
passing tests than Week 6's baseline, and the same two \texttt{xfail}
results and one \texttt{FAIL} as before, unrelated to this week's
additions. Neither \path{demos/coverage_reports/} nor
\path{demos/coverage_vs_pytest_cov/} is part of that run
(\texttt{pytest.ini} sets \texttt{testpaths = tests}); run each from
inside its folder.

% ============================================================
\section{How This Connects to the Reading}
% ============================================================
\begin{keyidea}
Weeks 5 and 6 defined coverage criteria. This week is about using the
tools: measure, read the gaps rather than the percentage, and turn each
gap into a behavior-level test with a real assertion. The demo folder
shows that method on the reading's own example. The two closed gaps
show it on code you have been working with since Week 4.
\end{keyidea}

% ============================================================
\section{Suggested Next Steps}
% ============================================================
\begin{enumerate}
\item From \path{demos/coverage_reports/}, reproduce the 74\% and 66\%
numbers, then open the HTML report (\texttt{coverage html}) and find
the three yellow lines (the reading's Exercise 3).
\item Write the two tests Exercise 4 asks for, then make the
\texttt{>=}~$\rightarrow$~\texttt{>} edit and confirm that one of them
fails. Restore the file afterwards.
\item Before reading the new tests in \path{tests/orders/} and
\path{tests/notifications/}, restore the Week 6 versions (see the demo
guide), look at the two gaps, and write your own tests for them. Then
compare.
\item Run \texttt{pytest -{}-cov=app -{}-cov-branch -{}-cov-report=html}
from the project root and use the Functions view to confirm that every
function in \path{app/} has been entered. Then name one function whose
100\% you would \emph{not} trust, and explain why.
\item In \path{demos/coverage_vs_pytest_cov/}, run the parallel-worker
pair (\texttt{pytest -n 2 -{}-cov} vs.\ \texttt{coverage run -m pytest
-n 2}) and explain the 67\% vs.\ 0\% difference in one sentence.
\end{enumerate}

\end{document}
